% LSPV-3 REPORT::SEEDBRANCH1::2026-10-05T00:43:00-07:00
\documentclass[11pt]{article}
\usepackage[margin=.88in]{geometry}
\usepackage{amsmath,amssymb}
\begin{document}
\section*{LSPV-3 SEEDBRANCH1: RMV direct-seed versus recursive-line accounting}
\textbf{STATUS: OPEN.} P versus NP remains open.
\section*{Source-level branch separation}
The vMS RMV definition has two distinct outputs at a node. First, for every MV line and every extractor seed it outputs the local-extractor string. Second, when $r>2$, it unions these strings with the sets returned by recursive calls on the line restrictions. Therefore extractor-seed enumeration is direct node work and is not a multiplicative factor on the recursive RMV child family.
\section*{Exact distinct-line quotient}
Let $B=\mathbb F_p^{r/2}$. An MV line in direction one is fixed by the second block, and an MV line in direction two is fixed by the first block. Hence there are exactly
\[
2|B|=2p^{r/2}
\]
distinct geometric MV lines. With explicit parent-table size $N=p^r$,
\[
\eta_{\rm line}=\frac12.
\]
Each line restriction has explicit table size $p^{r/2}=N^{1/2}$, so
\[
\theta_{\rm table}=\frac12.
\]
Canonical line identification from an anchor and direction consists of discarding the free block and retaining the fixed block; in the explicit coordinate representation this is polylogarithmic in $N$ and does not require SAT. This establishes the geometric quotient independently of semantic setup sharing.
\section*{Direct extractor branch}
The local extractor used by vMS has seed length
\[
s=O(r\log p)
\]
and output length $m$ after prefixing. Define the source-dependent constant
\[
\lambda:=\frac{s}{r\log_2p}.
\]
Then seed enumeration contributes
\[
2^s=N^{\lambda}
\]
up to lower-order encoding effects. The vMS lemma does not expose the hidden constant in its big-O statement, so no numerical value of $\lambda$ is source-certified here.
If one direct extractor output costs $N^{g_{\rm ext}+o(1)}$ in the chosen representation, direct work over all distinct lines and seeds has conservative exponent
\[
D_{\rm ext}=\frac12+\lambda+g_{\rm ext}.
\]
The line factor is required because the direct extractor is evaluated for every line as well as every seed.
\section*{Recursive branch}
The recursive child family has exponent
\[
D_{\rm rec}=\frac12+\frac12 C_{j}
\]
in the explicit-table normalization. Seed multiplicity does not enter this recursive term.
Thus the correct source-shaped recurrence is
\[
\boxed{
C_i\ge
\max\left\{
r_i,\;
\frac12+\lambda+g_{\rm ext},\;
\frac12+\frac12C_j
\right\}.
}
\]
For a self-similar recursive type, the recursive branch has least fixed point $1$. Therefore
\[
\boxed{
C=\max\left\{1,\frac12+\lambda+g_{\rm ext},r\right\}
}
\]
under the declared representation assumptions.
\section*{Correction to the previous heuristic}
The earlier candidate $\max\{\lambda+g_{\rm ext},\frac12+\frac12C\}$ omitted the fact that the extractor is run on every distinct line. The correct direct branch contains the additional line-enumeration exponent $1/2$.
Conversely, the earlier nested expression $\lambda+\frac12+\frac12C$ incorrectly multiplies seed enumeration into recursive calls. The source definition shows that this is also too pessimistic.
\section*{Threshold}
Ignoring node-local exponent $r$ for the moment, the source-side explicit-table recurrence is subquadratic exactly when
\[
\frac12+\lambda+g_{\rm ext}<2,
\]
i.e.
\[
\boxed{\lambda+g_{\rm ext}<\frac32.}
\]
It is at most linear only if
\[
\lambda+g_{\rm ext}\le\frac12.
\]
These are accounting thresholds, not SAT bounds.
\section*{SU05 coefficient audit}
SU05 gives explicit seed formulas for particular extractor parameterizations and its Hadamard binary conversion. However vMS Lemma 26 imports a local Reed--Muller extractor only at the theorem level with $s=O(r\log p)$. Without a source proof identifying the exact SU05 parameterization used in Lemma 26, importing a numerical leading constant would be an unjustified implementation choice. The coefficient therefore remains symbolic.
\section*{Space}
A materialized line child costs $N^{1/2+o(1)}$ space. Streaming/implicit evaluation can avoid simultaneous child materialization, but then generator workspace and recomputation are charged. Enumerating seeds needs only seed-counter space polylogarithmic in $N$ if outputs are consumed online; retaining all direct outputs would instead be polynomial space determined by their total output volume.
\section*{Acquisition and barriers}
This is source-side generator accounting only. It does not force an arbitrary polynomial-time SAT decider to construct RMV, expose line restrictions, or pay these costs. The result relativizes and is algebraic. The SAT acquisition/return theorem, SAME_REALIZATION bridge, forbidden-region comparison, nonrelativizing hinge and nonalgebrizing hinge remain OPEN.
\section*{Next target}
Derive $g_{\rm ext}$ for the declared SU05/Hadamard orbit realization under two explicit representations: random-access evaluation table and coefficient/LDE representation. Then determine whether an exact source-sanctioned seed parameterization yields a certified numerical $\lambda$, without conflating a permitted SU05 implementation with the theorem-level vMS extractor.
\section*{Status}
PROVED/EXACT for the source-shaped accounting: distinct-line exponent $1/2$, child-table exponent $1/2$, and separation of direct seed enumeration from recursive child multiplicity. CONDITIONAL: explicit-table exponent normalization and any $g_{\rm ext}$ value. OPEN: numerical $\lambda$ for the exact vMS imported extractor, SAT acquisition, and P versus NP.
\begin{thebibliography}{2}
\bibitem{SU05} R. Shaltiel and C. Umans, \emph{Simple Extractors for All Min-Entropies and a New Pseudorandom Generator}, J. ACM 52(2):172--216, 2005.
\bibitem{vMS25} D. van Melkebeek and N. Mocelin Sdroievski, \emph{Instance-Wise Hardness and Refutation versus Derandomization for Arthur--Merlin Protocols}, Computational Complexity 34:15, 2025.
\end{thebibliography}
% END_OF_VERIFIED_REPORT::LSPV-3::SEEDBRANCH1
\end{document}
